\documentclass{article}
\usepackage[utf8]{inputenc}
\usepackage[spanish]{babel}

\title{\textbf{Proyecto Final - Machine Learning} \\[0.5em] \Large Predicción del exito académico en educación superior}
\author{Javier Martínez Berenguer}
\date{\today}

\begin{document}

\maketitle

\section{Abstract}


\section{Preprocesado de los datos}

En esta sección se detallan los procedimientos llevados a cabo en la explración y el 
tratado del conjunto de datos para garantizar 
la calidad de los mismos y la integridad de los modelos empleados. 
El dataset original consta de 4,424 registros y 37 variables (19 numéricas y 18 categóricas).

\subsection{Exploración inicial y análisis descriptivo (EDA)}

\subsubsection{Distribución de la variable objetivo y desbalanceo}
La variable \texttt{objetivo} presenta una distribución asimétrica: \textit{Graduado} (2,209 instancias, 50.0\%), \textit{Abandono} (1,421 instancias, 32.1\%) y \textit{Matriculado} (794 instancias, 17.9\%). 
El desbalanceo de la clase \textit{Matriculado} 
    sugiere que el \textit{Accuracy} no será una métrica fiable. Un modelo ingenuo podría ignorar esta 
    clase y aun así obtener un rendimiento aparente alto. Es por ello que se llevará a cabo el análisis de la matriz de 
    confusión y el F1-Score por clase.


\subsubsection{Análisis de las variables académicas}
A partir del análisis de las distribuciones, se observa que las variables de desempeño del primer semestre parecen ser los predictores más influyentes. Esto se debe a que muestran una clara diferenciación entre los grupos de alumnos, especialmente en el caso de las notas medias, donde aquellos valores bajos están fuertemente ligados al abandono.
\begin{itemize}
    \item \textbf{Efecto de las notas bajas:} Al observar la \texttt{nota\_media\_1sem}, se distinguen dos grupos muy claros: la mayoría de los alumnos tiene notas entre 10 y 15, pero hay un grupo numeroso con un 0.0. Este grupo está compuesto casi totalmente por alumnos que terminan en \textit{Abandono}, lo que indica que dejaron de asistir incluso antes de la primera evaluación.
    \item \textbf{Gestión de la multicolinealidad:} En el bloque de variables académicas (asignaturas matriculadas, evaluadas y aprobadas) existe una correlación superior a 0,85. Debido a la alta correlación, se han adoptado dos estrategias en función del modelo a emplear para la clasificación.
\begin{itemize}
    \item \textbf{Regresión Logística:} Se aplica regularización Lasso ($L_1$). Es necesaria porque la redundancia infla la varianza de los coeficientes, volviendo el modelo inestable. Lasso resuelve esto forzando a cero los coeficientes de las variables menos informativas.
    \item \textbf{Random Forest:} Es robusto ante la colinealidad. Al utilizar \textit{bagging} (selección aleatoria de un subconjunto de variables en cada nodo), el modelo evita depender de una única variable correlacionada.
\end{itemize}
\end{itemize}

\subsubsection{Factores socioeconómicos y demográficos}
El análisis revela patrones críticos fuera del ámbito académico:
\begin{itemize}
    \item \textbf{Becados y deudores:} Los estudiantes con la variable \texttt{becado} en 'si' muestran una tasa de graduación significativamente mayor, mientras que los \textit{deudores} presentan una probabilidad de abandono tres veces superior a la media. Esto coloca a las variables económicas como predictores con un impacto considerable
    \item \textbf{Sesgo de Edad:} La distribución de \texttt{edad\_al\_matricularse} presenta un sesgo positivo pronunciado (asimetría de 2.05). Los estudiantes más mayores (outliers por encima de los 34 años) muestran trayectorias académicas más erráticas, posiblemente debido a la dificultad de conciliar la vida laboral y académica.
    \item \textbf{Variables macroeconómicas:} Variables como la \texttt{tasa\_desempleo} o el \texttt{PIB} presentan una varianza baja dentro del dataset. Su análisis descriptivo indica que su correlación individual con el éxito académico es débil, lo que plantea su posible exclusión para simplificar el modelo sin pérdida de precisión.
\end{itemize}

\subsection{Valores faltantes}
El dataset no contiene valores nulos técnicos 
(\texttt{NaN}), sino valores nulos codificados como etiquetas de texto. Dichos valores se presetan
únicamente en valores pertenecientes a variables categóricas. En variables como \textit{ocupación\_madre} 
y \texttt{cualificación\_padre}, se detectaron términos como \textit{'desconocido'}, \textit{'en\_blanco'}
y \textit{'otra\_situación'}.

\begin{itemize}
    \item \textbf{Decisión:} Se han unificado estas etiquetas en una categoría única 
    denominada \texttt{'informacion\_no\_disponible'}.
\item \textbf{Justificación:} Simplifica las variables y evita el sobreajuste 
(\textit{overfitting}). Al agruparlas, se refuerza la señal de que la falta de 
información familiar es un indicador clave de riesgo académico.
\end{itemize}

\subsection{Valores atípicos}
Para la detección de valores atípicos en variables numéricas se ha empleado el método del Rango 
Intercuartílico ($IQR$).

\begin{itemize}
    \item \textbf{Variables académicas:} Se detectaron valores extremos poco realistas, como registros con 45 asignaturas evaluadas en un solo semestre. Se aplicó \textit{capping} para limitar estos casos a los umbrales del $IQR$, evitando así que estos datos aislados distorsionen el aprendizaje del modelo.
    \item \textbf{Edad:} Se han conservado los valores atípicos en \texttt{edad\_al\_matricularse} (ej. alumnos mayores de 34 años), ya que representan perfiles demográficos reales y necesarios para el estudio.

\item \textbf{Calificaciones (Notas medias):} El tratamiento de los valores de 0.0 se diferencia según la tarea:
\begin{itemize}
    \item \textbf{Clasificación y No Supervisado:} Se mantienen. Estos valores son un indicador clave del perfil de \textit{Abandono}. Eliminarlos reduciría drásticamente el F1-Score y dificultaría la creación de clústeres de estudiantes en riesgo.
    \item \textbf{Regresión:} Se filtran estos valores. Al reflejar una deserción prematura y no la capacidad académica real del alumno, su inclusión distorsionaría la predicción de la nota final, sesgando los resultados del modelo.
\end{itemize}
\end{itemize}

\end{document} 